\appendix
\section{Expanded dessin, resolution and derived calculations}
\subsection{The cycle order in coordinates}
Number a permitted cycle in travel order $0,\ldots,\ell-1$. Map an arrow $i\to i+1$ to $E_{i,i+1}$ for $i<\ell-1$, and the final arrow to $u_c E_{\ell-1,0}$. A path from $i$ to $j$ then gives $u_c^rE_{ij}$, where $r$ counts cut crossings. These are linearly independent path normal forms, and every positive-length permitted path has exactly one such expression on exactly one branch. The full-cycle sum in a branch is $u_c I$. At a paired edge its two diagonal occurrences come from a single original idempotent; no positive-length coefficients are identified. This proves the order description including self-pairings within one cycle. Although the suborder is generically a product of full matrix algebras, it is not their unglued product at the special fiber.

For the interval this becomes the literal fiber product
\[
 \Lambda=\{(f,g)\in k[[X]]^2:f(0)=g(0)\},\qquad
 x=(X,0),\ y=(0,X),\ u=(X,X).
\]
Every element has a unique expression $a(u)+b(u)x$: take $a=g$ and $b=(f-g)/X$. Thus it is $R$-free on $1,x$ with $x^2=ux$ and $y=u-x$. The annihilator of $x$ consists of pairs $(0,g)$ with $g(0)=0$, exactly $y\Lambda$; the other branch is symmetric. Multiplication by $u$ is injective, coefficientwise. These are actual formal power-series statements, not conclusions from finite truncations.

\subsection{The two-arm complex and its calculations}
For an arrow $a$ write $f(a)$ for its unique forbidden successor. The summand in $P_n(i)$ indexed by the initial arrow $a$ has endpoint $t(f^{n-1}(a))$. Its map multiplies by that last arrow. Thus the resolution can be computed by iterating a finite permutation, while the projectives themselves remain infinite-dimensional. The Hom complex against a module $N$ is the direct sum of the endpoint spaces with differentials given by the same arrow matrices. Tensor against a left module instead uses the same path entries acting on the left. No rank or exactness claim for an arbitrary $N$ is inferred by counting endpoints alone.

For the nodal residue field the maps are particularly explicit:
\[
 \cdots\xrightarrow{\operatorname{diag}(x,y)}\Lambda^2
 \xrightarrow{\operatorname{diag}(y,x)}\Lambda^2
 \xrightarrow{(x\ \ y)}\Lambda\longrightarrow k\longrightarrow0.
\]
If $xz_1+yz_2=0$, its first branch says the first branch of $z_1$ is zero, and its second branch says the second branch of $z_2$ is zero. Hence $z_1=yw_1$, $z_2=xw_2$. The same annihilator calculation proves the alternating kernels in all degrees. Projectivity is genuine: the terms are regular modules or their finite direct sums. It is not a finite free resolution, despite the period of its positive tail.

For the elliptic theta the quiver vertices are $v_0=\{0,3\}$, $v_1=\{1,4\}$, $v_2=\{2,5\}$. The arrows $a_0,a_3$ go $v_0\to v_2$, $a_2,a_5$ go $v_2\to v_1$, and $a_1,a_4$ go $v_1\to v_0$. The two forbidden arms from $v_0$ are
\[
 a_0,a_5,a_1,a_3,a_2,a_4,\ldots,\qquad
 a_3,a_2,a_4,a_0,a_5,a_1,\ldots.
\]
Their top pairs are $(v_2,v_2),(v_1,v_1),(v_0,v_0)$, repeated. The three steps exchange the arms. Their distinct top pairs exclude shorter unordered positive periods. The spherical theta instead has forbidden two-cycles and successive top pairs $(v_2,v_1),(v_0,v_0)$, giving period two. Thus even identical hereditary overorders need not give identical resolutions.

Quadratic Koszul duality here has a directly checkable local meaning: relations complement the two possible successors. The dual dessin replaces the graph rotation by $\phi^{-1}$. In the Ext algebra against $A/J$, concatenating the forbidden-word basis gives the product when the join is forbidden; a permitted join gives zero. One can see this by lifting a length-$r$ basis class to the resolution: discard the prescribed initial $r$ letters and identify the remaining summand with its suffix resolution, taking zero on other prefixes. These prefix-discard comparison maps commute with the suffix differentials, with the standard shift signs when viewed as shifted complexes. The concatenated indexing word exists exactly when its join is forbidden; evaluation of the two basis classes on that word gives one. If the join is permitted there is no such indexing word and the product is zero. Composition therefore concatenates indexing words, with the chosen right-module/Yoneda opposite convention. This is a graded algebra statement; it neither makes the unbounded resolution perfect nor constructs an equivalence of all formal-completion derived categories.

For an arbitrary bounded projective complex $P^\bullet$, its determinant line over a commutative coefficient ring is $\bigotimes_n(\det P^n)^{(-1)^n}$, with determinant duals for negative exponents. A chain automorphism acts on this line; a quasi-isomorphism gives the corresponding determinant-functor isomorphism. For a general endomorphism $F$, the expression for $1-zF$ uses the locus where it is invertible and gives a rational determinant-line action/trivialization, not an everywhere invertible line map at singular parameters. Additivity is compatible with exact triangles under the usual determinant-functor conventions \cite{KnudsenMumford,Breuning}. In our order context an $R$-linear determinant uses the finite free underlying $R$-modules; it must specify a separate endomorphism. Determinants of a differential presentation, alternating determinants of a chain action, regularized analytic determinants and the Cartan path matrix are four distinct constructions. Their equality requires a proof, not their shared name.

The geometric model literature \cite{GeometricGentle,GeometricResolutions,GeometricKoszul} is useful for identifying which crossings and gradings would have to be represented by a functor. To apply it here one must specify the category of polynomial strings, completed lattices or finite-length quotients; construct its surface/curve objects; and prove that degree-$r$ intersections compute the Hom group into shift $r$, with local systems and boundary exceptions included. We have supplied explicit walks and maps in the examples, not claimed this general completed-order theorem. On a finite word, a map can already be computed rigorously without such a theorem: assign a coefficient matrix between equal-vertex occurrences and solve $F_{t(a)}N(a)=N'(a)F_{s(a)}$ for every arrow. The diamond calculations below implement exactly this system.

\section{Semi-invariants, tableaux and explicit intertwiners}
\subsection{Derived information retained by a presentation}
Let $K=\operatorname{im}f$. The injection $\operatorname{Hom}(K,N)\to\operatorname{Hom}(P_1,N)$ descends to the Ext injection used in Theorem~\ref{thm:semi}; its kernel modulo $\operatorname{im}\operatorname{Hom}(P_0,N)$ is zero because the map $P_1\to K$ is onto. The missing part of the cokernel detects maps on $P_1$ not descending through $K$. This explains exactly why an injective resolution map is needed for its full Euler interpretation.

In the node example the resolution of $M=\Lambda/(x)$ has maps $x,y,x,y,\ldots$. On $N=k[y]/(y^2)$ the Hom cochain maps are $0,y,0,y,\ldots$. Thus $\operatorname{Hom}(M,N)=N$ and $\operatorname{Ext}^1(M,N)=\ker y$ have dimensions two and one; the cokernel of the first zero map is two-dimensional. In the elliptic-square example the first four Hom maps are $1,0,0,0$, periodically. Hence degrees zero and one vanish, degrees two and three are $k$, and degree four vanishes. This also makes clear why a nonzero scalar semi-invariant does not give a derived orthogonality theorem.

With coordinate action $(gF)(N)=F(g^{-1}N)$ the Schofield weight is $a-b$, the negative of the point-evaluation weight. For a single arrow it is the same elementary distinction as the action on $\operatorname{Hom}(V,W)$ versus on its coordinate functions $V\otimes W^*$. The Cauchy decomposition for coordinate functions is
\[
 \operatorname{Sym}^m(V\otimes W^*)\cong
 \bigoplus_{\lambda\vdash m}S_\lambda(V)\otimes S_\lambda(W)^*.
\]
A relation quotient can kill these components; its ambient multiplicities do not descend unchanged. Acyclic-quiver semi-invariant/LR results such as \cite{DerksenWeyman} retain their acyclicity assumptions. The explicit two-arrow slice below is acyclic, whereas the whole theta quiver is not.

\subsection{One Young operator and one LR decomposition in coordinates}
For $\lambda=(2,1)$ number the first row $1,2$ and second row $3$. Then
\[
 Y=(I+P_{12})(I-P_{13}),\quad Y^2=3Y,\quad
 \dim\operatorname{im}(Y\mid(k^2)^{\otimes3})=2.
\]
The hook product is three; its Specht dimension is $3!/3=2$, also two here for a different reason. The normalized projector extracts one $S_{(2,1)}$ copy. The trace $\tr((Y/3)M^{\otimes3})$ is therefore the Schur character, not twice that character. This fixes the common row/column normalization ambiguity.

Write $t=\tr M$, $\Delta=\det M$ for a two-dimensional endomorphism. Its complete characters obey $h_0=1$, $h_1=t$, $h_m=th_{m-1}-\Delta h_{m-2}$, and
\[
 \chi_{(a,b)}(M)=\Delta^b h_{a-b}(M),\quad a\ge b\ge0.
\]
In particular $h_2=t^2-\Delta$, $h_3=t^3-2t\Delta$, and $\chi_{(2,1)}=t\Delta$. Thus $t h_2=h_3+\chi_{(2,1)}$ is the LR identity $V\otimes\operatorname{Sym}^2V\cong\operatorname{Sym}^3V\oplus S_{(2,1)}V$, with $S_{(2,1)}V\cong\det V\otimes V$.

An explicit natural splitting avoids relying just on characters. In characteristic zero multiplication $m:V\otimes\operatorname{Sym}^2V\to\operatorname{Sym}^3V$ has natural section
\[
 s(v_1v_2v_3)=\tfrac13\sum_{i=1}^3v_i\otimes\prod_{j\ne i}v_j.
\]
The map $j:\bigwedge^2V\otimes V\to V\otimes\operatorname{Sym}^2V$ sends $(v\wedge w)\otimes z$ to $v\otimes wz-w\otimes vz$. In dimension two it is injective: for a basis $e_1,e_2$ its images at $z=e_1,e_2$ have distinct independent coordinates. They lie in $\ker m$, whose dimension is $6-4=2$. Hence $(s,j)$ is the required isomorphism. Tensorial definitions make both maps commute with every linear endomorphism, including singular arrow maps. This is a genuine LR intertwiner on our dessin representation stratum.

For a concrete theta representation take
\[
 X_1=\begin{pmatrix}1&1\\0&1\end{pmatrix},\quad
 X_2=\begin{pmatrix}2&0\\0&1\end{pmatrix},\quad
 X_3=\begin{pmatrix}1&0\\1&1\end{pmatrix}.
\]
Then $M=X_3X_2X_1=\left(\begin{smallmatrix}2&2\\2&3\end{smallmatrix}\right)$, $t=5$, $\Delta=2$, $h_2=23$, $h_3=105$ and $\chi_{(2,1)}=10$. The identity is $115=105+10$. Since the other permitted branch is zero, every forbidden product is zero; this is an actual representation of the infinite elliptic dessin algebra.

In the invertible-arrow open stratum one can gauge $X_1,X_2$ to the identity, leaving conjugation of $M$. The invariant-ring proof in the core works also at singular arrows by cycle-trace generation, not by assuming that this gauge choice extends there. The degree-one monodromy coefficients have total arrow degree three; this is why the Hilbert-series denominator uses $3j$, not $j$.

For the Cauchy kernel choose the Schur subspace inside the tensor power with its inherited Hilbert norm. $M^{\otimes m}$ preserves it and has norm at most $\|M\|^m$. This proves the trace bound for nonnormal matrices as well as diagonal ones. Formal Cauchy identities extend by polynomial continuity in matrix entries; the stated absolute bound then justifies analytic summation. Gram positivity does not assert entrywise positivity or positivity of an unrelated arithmetic form. For three exact test matrices $M$, $\left(\begin{smallmatrix}1&1\\0&2\end{smallmatrix}\right)$, $\left(\begin{smallmatrix}0&1\\-1&1\end{smallmatrix}\right)$, the choice $r=1/100$ satisfies $rR_0^2\le21/100<1$ using Frobenius norms, and all seven principal minors of the rational kernel matrix are nonnegative.

\subsection{Two precise functorial limitations}
At any actual vertex simple $S$ of $A$, the split exact sequence $0\to S\to S\oplus S\to S\to0$ uses first inclusion and second projection. Termwise $\operatorname{Sym}^2$ has dimensions $1,3,1$ at that vertex; its middle kernel contains the mixed monomial as well as the first square, but its first image contains only the first square. Middle homology is $S$. Thus even split exactness is not preserved. Genuine derived Schur constructions instead use mixed tensor degrees, signed total differentials and suitable flat/projective replacements \cite{SchurComplexes,DerivedSchur}; their membership in a particular order category must still be checked.

The $R$-lattice $M=R e_x\oplus R e_y$ with $x=\operatorname{diag}(u,0)$, $y=\operatorname{diag}(0,u)$ has $x+y=uI$. On $\operatorname{Sym}^2_RM$ the images are $\operatorname{diag}(u^2,0,0)$ and $\operatorname{diag}(0,0,u^2)$. Their sum kills the mixed vector, hence equals neither $uI$ nor $u^2I$. This is a specific failure to retain the original central action, not a failure of the polynomial monomial representation functor. Even a common ramified base substitution does not repair this mixed vector.

Semi-invariants retain useful moduli on nilpotent modules. In the theta quiver keep only $a_0,a_3:v_0\to v_2$, put their scalar values $(x,y)=(1,\lambda)$ and all other arrows zero. The $v_1$ simple splits off. Every length-two path acts by zero, so these are continuous completed modules. The two scalars have the same weight $g_{v_2}/g_{v_0}$, and $y/x$ is a rational invariant on $x\ne0$, giving the projective pencil $[x:y]$. It is not a regular invariant on the whole affine locus, whose orbit closures meet zero.

For the interval quotient take $x=J_2$, $y=\lambda J_2$. With its top and socle specified, the separated Kronecker coordinates are again $(1,\lambda)$; the relation to the same-weight ratio is exact. The pp formula in the core is stronger than merely noticing a ratio: on $M_\lambda$ its solution set is $\operatorname{im}x$, whereas on $M_\mu$, $\mu\ne\lambda$, the equation $(\mu-\lambda)xv=0$ forces $z=xv=0$. Different parameters are therefore distinguishable by actual positive-primitive conditions even though every cycle trace is zero.

\section{Full lattice and diamond diagrams}
\subsection{The nodal exact diagram}
The normalization proof of Theorem~\ref{thm:nodear} uses $BM$ inside the generic module. Replacing it without justification by $B\otimes_\Lambda M$ could introduce torsion and would not prove the same statement. Once the subdirect residue subspace is decomposed, the branch basis changes lift over complete DVRs; taking inverse images modulo the conductor gives the direct decomposition, not just a decomposition of residues.

Here is the promised diagram, with the terminal short exact column for every position:
\[
\begin{tikzcd}[column sep=large,row sep=large]
0\arrow[r]&Y\arrow[r,hook]\arrow[d,"u"]&\Lambda\arrow[r,"x"]\arrow[d,"u"]&X\arrow[r]\arrow[d,"u"]&0\\
0\arrow[r]&Y\arrow[r,hook]\arrow[d]&\Lambda\arrow[r,"x"]\arrow[d]&X\arrow[r]\arrow[d]&0\\
0\arrow[r]&Y/uY\arrow[r,hook]&\Lambda/u\Lambda\arrow[r]&X/uX\arrow[r]&0
\end{tikzcd}
\]
Each vertical column is completed by zero above and below. In the bottom row identify $Y/uY=k$ via generator $y$, $X/uX=k$ via $x$, and $\Lambda/u\Lambda=k[\epsilon]/(\epsilon^2)$ with $x=\epsilon$, $y=-\epsilon$. Its first map is $1\mapsto-\epsilon$, its second the augmentation. It is nonsplit as a dual-number sequence. This explicit order reduction, not an unrelated finite algebra, explains its appearance.

For a further Ext calculation, resolve $X\cong\Lambda/(y)$ with successive maps $y,x,\ldots$. On $Y$ the first is multiplication by $y$ and the second zero, so $\operatorname{Ext}^1(X,Y)=Y/yY=k$; on $X$ the first is zero and the second injective, so $\operatorname{Ext}^1(X,X)=0$. Interchanging branches gives the other two cases. The two nonsplit lattice rows realize the nonzero cross classes. Their Hom spaces are finite over $R$, generally infinite over $k$.

\subsection{Chebyshev geometry and the exact quotient}
The identity $1-T_n(z)^2=(1-z^2)U_{n-1}(z)^2$ follows from $z=\cos\theta$ and the sine/cosine multiple-angle formulas, hence as a polynomial identity. $T_n'=nU_{n-1}$. Therefore $\beta_n$ has its only zero at infinity, of order $2n$; its $n$ points over one are double since $\beta_n-1=T_n^2/(1-T_n^2)$; and its poles are simple at $\pm1$ and double at the $n-1$ interior roots of $U_{n-1}$. The derivative has no additional critical values. The total ramification is $(2n-1)+n+(n-1)=4n-2$, as required for a degree-$2n$ rational Belyi map. The dual graph has $n+1$ vertices of degrees $1,1,2,\ldots,2$, and $n$ edges; it is a connected path. Its planar ribbon embedding is unique, so this identifies the particular nested-loop dessin, not only its passport.

In the quotient retaining orientation $\Omega$, any surviving path traverses each chain edge at most once and either endpoint loop at most once. Its length is at most $n+1$; thus the killed-arrow ideal contains $J^{n+2}$. The quotient of $A_n$ and of its completion is the same $H_n(\Omega)$. For $n=3$ with forward chain arrows, its path counts in lengths zero through four are $3,4,3,2,1$, total thirteen. These counts describe a quotient representation category of the infinite order, not a change of primary object.

Inflation is full because the two algebras have exactly the same actions on these annihilated objects. It preserves internal exact sequences. It does not preserve arbitrary ambient extension closure: the two adjacent simples are annihilated by the killed reverse arrow, but the two-dimensional extension having only that arrow nonzero is not. Likewise endpoint-local freeness over $k[l]/l^2$ or $k[r]/r^2$ in \cite{HuangLinSu} is not $R$-freeness. To transfer an almost-split property to a larger category one must prove the requisite approximation/factorization statements; locally finite quiver criteria \cite{BautistaLiuPaquette,Paquette,LiuNgPaquette} do not apply just because each vertex has two arrows. The permitted cycles violate interval/Hom finiteness.

\subsection{Nine actual strings and all maps}
Take $n=3$ with $1\xrightarrow{a}2\xrightarrow{b}3$, loops $l,r$ and $l^2=r^2=0$ in the quotient. Here $b$ denotes the retained forward $a_2$, not a reverse arrow of the original doubled chain. The following words are in travel order. A superscript $-$ means a formal inverse in a word, not an inverse algebra element.
\begin{center}\small
\begin{tabular}{llll}
\toprule
Object & word & start & dimension vector\\\midrule
$A$&$l^- a b r b^- a^- l^-$&1&$(4,2,2)$\\
$B$&$l^- a b$&1&$(2,1,1)$\\
$C$&$l^- a b r^-$&1&$(2,1,2)$\\
$D$&$a b r b^- a^- l^-$&1&$(3,2,2)$\\
$E$&$a b$&1&$(1,1,1)$\\
$F$&$a b r^-$&1&$(1,1,2)$\\
$G$&$r^- b^- a^- l a b r b^- a^- l^-$&3&$(4,3,4)$\\
$H$&$r^- b^- a^- l a b$&3&$(2,2,3)$\\
$I$&$r^- b^- a^- l a b r^-$&3&$(2,2,4)$\\\bottomrule
\end{tabular}
\end{center}
A basis vector is assigned to each visited position $0,\ldots,\ell(w)$ at its vertex. A direct letter sends its preceding vector to its following vector; an inverse letter sends the following vector to the preceding vector. Sum these elementary maps for each arrow and put all unmentioned entries zero. This defines every matrix, including repeated loops, and verifies the relations. Set the killed reverse arrows to zero to obtain actual $\Lambda_3$-modules.

Their signed grid is
\[
\begin{tikzcd}[column sep=large,row sep=large]
 A\arrow[r,"p"]\arrow[d,"-p"']&B\arrow[r,"i"]\arrow[d,"p"]&C\arrow[d,"-p"]\\
 D\arrow[r,"p"]\arrow[d,"-i"']&E\arrow[r,"i"]\arrow[d,"i"]&F\arrow[d,"-i"]\\
 G\arrow[r,"p"]&H\arrow[r,"i"]&I.
\end{tikzcd}
\]
Horizontal projections keep the initial word positions and kill the remaining positions. Horizontal inclusions append the final $r^-$, keeping old positions. Upper vertical projections delete the initial $l^-$: old position $j\ge1$ maps to position $j-1$, and old position zero maps to zero. Lower vertical inclusions prepend $r^-b^-a^-l$: position $j$ maps to $j+4$. Multiply these maps by the indicated signs. For every arrow these formulas commute with the string action, including the two positions adjacent to a cut or join; those are exactly where the loop inverse and chosen orientation matter.

Every small square has two length-two composites whose sum is zero. The exact sequences are
\[
\begin{split}
0&\to A\to B\oplus D\to E\to0,\qquad
0\to B\to C\oplus E\to F\to0,\\
0&\to D\to E\oplus G\to H\to0,\qquad
0\to E\to F\oplus H\to I\to0.
\end{split}
\]
For example the maps out of $A$ retain its initial positions $0,1,2,3$ in $B$ and delete its first position in $D$, so their joint kernel is zero. The maps to $E$ are onto, and the dimension vectors give $\dim A+\dim E=\dim B+\dim D$ vertexwise. The zero composite plus this equality proves exactness. The other three use precisely the same keep/delete or append/prepend rules; their vertexwise dimension equalities and jointly injective/surjective maps prove exactness. The shipped exact checker constructs all matrices and verifies these assertions independently over $\mathbb Q$.

The stronger almost-split property is imported by identifying $E=w=ab$ as the minimal string of type $(2,2)$ and these exact hook/cohook words with \cite[Proposition3.12]{HuangLinSu}, using that paper's Propositions2.6/2.9 and Lemma3.10. Opposite-module conventions are reconciled by reversing travel/composition conventions; the closed doubled quiver is isomorphic to its opposite by exchanging the two arrows at each chain edge. Thus intrinsically $\tau E=A$, $\tau F=B$, $\tau H=D$, $\tau I=E$. The rows and columns of this grid are \emph{not} generally short exact, unlike the nodal diagram. Two small cautions in that source are not needed here: a tube-rank label in Example3.9(3) is inconsistent with its listed mouth, and a general homogeneous-tube statement must not be read as a converse classifying every such object as a band. Neither is used to import these four meshes.

\subsection{The pure-injective quotient and what its dimension means}
For $B_0=\Lambda/J^2=k\langle x,y\rangle/(x,y)^2$, words alternate direct and inverse letters because every direct/direct and inverse/inverse adjacent pair is forbidden. Up to rotation and inverse its only primitive band is $xy^{-1}$; the bridge quiver has two oriented band vertices with no connecting nontrivial bridge. The longest path has length zero. The domestic string results \cite{PuninskiPrest,LakingPrestPuninski} therefore give relative dimension $0+2=2$, not an automatic dimension three for every one-domestic algebra. Finite-dimensional points have relative CB rank zero; the appropriate Prüfer/adic and one-sided periodic string points have rank one; the unique generic point has rank two. These statements concern the quotient category with its inherited pp structure.

Indeed a pp formula is a finite system of linear equations with existentially quantified unknowns. On a module annihilated by $J^2$ every coefficient reduces to $B_0$, so purity and pure-injectivity are preserved in both directions under inflation. The condition $MJ^2=0$ is a definable closed condition, giving the closed Ziegler subspace. Its relative ranks need not be ranks in the entire $\Lambda$ spectrum. This subcategory is not extension closed: $\Lambda/(x^3,y)$ gives an extension of objects killed by $J^2$ that is not itself killed by $J^2$. Since $u^2$ annihilates every object of the subcategory, it contains no nonzero $R$-free lattice. Ringel's representation-theoretic node construction is also not a Morita equivalence identifying the geometric nodal order with the Kronecker algebra. Infinite-dimensional polynomial string classifications \cite{CrawleyBoeveyStrings} retain finitely-controlled or pointwise-artinian assumptions; an arbitrary complete lattice does not automatically satisfy them as an algebraic $A$-module.

\section{The operator extension in detail}
\subsection{A corrected cycle lift and the integer boundary map}
In the permitted-state graph each length-$\ell$ cycle has a full matrix-corner Laurent/circle coordinate. Its chosen $K_1$ generator is the full cycle at one state plus the complementary identity. In Fock space $L_c^*L_c=Q_a+v_{s(a)}$. Thus the tempting lift $L_c+(1-Q_a)$ has overlapping initial pieces and is not an isometry. The corrected lift in Theorem~\ref{thm:fock} first restricts to $Q_a$:
\[
 W=L_cQ_a+(1-Q_a),\quad W^*W=I,\quad
 I-WW^*=Q_a-L_cQ_aL_c^*.
\]
A path starting with $a$ and length greater than $\ell$ has uniquely the full prefix $c$ followed by another path starting with $a$. The missing range consists exactly of lengths $1,\ldots,\ell$. Its terminal multiset is the original graph's incident-edge multiset at that cycle vertex. These finite defect vectors all lie in terminal compact ideals; their classes give the claimed unsigned incidence entries. The sign convention affects only the overall sign of the boundary map.

Here is the Smith calculation without spectral assumptions. Choose a spanning tree of the connected underlying graph and one root. By changing signs on one bipartition of the tree, its $E\times V$ incidence submatrix becomes the ordinary oriented tree incidence, with $V-1$ unit pivots. After eliminating these pivots, a remaining column represents the common root coordinate, and each nontree edge evaluates on it as zero when its fundamental cycle is even, or $\pm2$ when odd. A loop is the latter case. If all are even the graph is bipartite and this column is zero; otherwise its entries have gcd two, giving the final pivot two. Remaining rows are free. Consequently the kernel/cokernel, then the six-term sequence, give the $K$ groups stated in the core. No theorem reconstructing a cyclic embedding from those groups is asserted.

The compact ideal is block diagonal, not $\mathcal K(\mathcal H)$: prefixes never move terminal vertices, so cross-terminal rank-one maps do not belong to the displayed ideal. This is why the arithmetic insertion must choose one terminal sector. A graph/Cuntz boundary quotient kills the vacua and hence this ideal, so that insertion need not survive passage to the quotient.

\subsection{Spectral domain and arithmetic trace}
A vector belongs to the adjoint domain of the diagonal $N$ precisely when its coefficientwise product by the real sequence $|p|$ is square summable; this proves self-adjointness on the given domain. The truncations by length prove finite supports are a core. Resolvent coefficients $(|p|-i)^{-1}$ tend to zero and each length space is finite dimensional, giving compactness. Prefix and deletion shift the length by at most one, so preserve the domain and prove the commutator formulas there. Leibniz then handles the dense polynomial *-algebra.

Every terminal sector has one vacuum and two paths at each length $n\ge1$, even when some graph endpoints coincide: distinct formal paths remain distinct basis vectors. Thus restriction to a sector still has a simple vacuum, but it need not be a faithful representation of the whole concrete algebra. The trace of $e^{-tN^2}$ is $E+2E\sum_{n\ge1}e^{-tn^2}$, distinct from the semigroup trace of $N$ in the core. The trace identity with $\chi(N)$ is only on $\Rea s>1$, where the product is trace class. Analytic continuation of the scalar $L$-function does not make that same operator product trace class on $\Rea s>0$.

\section{Abel determinant estimates and the arithmetic program}
For completeness the inserted family has a full elementary convergence proof. Since complete nonprincipal character periods sum to zero, $|A_n|\le B_0\le q_\chi$. On $2^k\le n<2^{k+1}$,
\[
 \sum |u_n|^2\le\frac{B_0^2}{(k+1)^2},\qquad
 \sum |v_n(s)|^2\le M^2(k+1)^2r^k,\quad r=2^{-2\sigma_0}<1.
\]
The second follows from $n^{-s}-(n+1)^{-s}=s\int_n^{n+1}x^{-s-1}\,dx$. Define
\[
 Q_J(r)=r^J\left(\frac{(J+1)^2}{1-r}
 +\frac{2(J+1)r}{(1-r)^2}+\frac{r(1+r)}{(1-r)^3}\right).
\]
Expanding $(J+j+1)^2$ in the three geometric moment sums proves $Q_J=\sum_{k\ge J}(k+1)^2r^k\to0$. Locally uniform $\ell^2$ convergence of the holomorphic finite-coordinate vectors proves Banach holomorphy. The linear functional $z\mapsto\sum v_nz_n$ has norm $\|v\|$; this is the complex-linear coefficient pairing, not a conjugate-linear dependence on $s$.

Write $\mathcal I(v)(z)=\sum_n v_nz_n$ for the coefficient functional, distinct from the arrow ideal $J$. For $N_J=2^J-1$, $J\ge1$, let $K_N=u\otimes(\mathcal I(P_Nv)-e_1^*)$. Then
\[
 \|K-K_{N_J}\|_1\le\sqrt2 B_0M\sqrt{Q_J(r)}.
\]
Finite summation by parts gives
\[
 T_N=\det_F(I+K_N)
 =\sum_{n\le N}A_n(n^{-s}-(n+1)^{-s})
 =S_N-A_N(N+1)^{-s}.
\]
Adding the last endpoint times $e_1^*$ to the functional recovers $S_N$ exactly because $u_1=1$. Its trace-norm error is at most
\[
 \sqrt2 B_0\bigl(M\sqrt{Q_J(r)}+B_0(N_J+1)^{-\sigma_0}\bigr).
\]
For arbitrary cutoffs use the last whole dyadic block for an enclosing tail. Galerkin compressions also converge: $\|u-P_{N_J}u\|\le B_0/\sqrt J$. With $C_K=1+M\sqrt{(1+r)/(1-r)^3}$,
\[
 \|K-P_{N_J}KP_{N_J}\|_1
 \le\frac{B_0 C_K}{\sqrt J}+\sqrt2 B_0M\sqrt{Q_J(r)}.
\]
The raw endpoint adds the same endpoint bound as before. Isometric insertion by $U$ preserves nonzero singular values, every norm estimate, and the Fredholm determinant. Everything is in one compact ideal of the actual infinite dessin Fock algebra.

For primitive parity $a$ use $g(s)=(q_\chi/\pi)^{(s+a)/2}\Gamma((s+a)/2)$, holomorphic nonzero on $\Rea s>0$. The completed functional is $g(s)\mathcal I(v)-e_1^*$, and the completed raw functional is $g(s)(\mathcal I(P_Nv)+A_N(N+1)^{-s}e_1^*)-e_1^*$. Their determinants are exactly $gL$ and $gS_N$. Scaling the entire old perturbation would instead give $1+g(L-1)$, the wrong normalization. Trace-class background is standard \cite{Bornemann}, while these special rank-one identities follow directly from $\det_F(I+u\otimes\varphi)=1+\varphi(u)$.

The metric obstruction is likewise elementary. With inner product conjugate-linear in its first variable, symmetry of $T=u\otimes\varphi$ implies
\[
 \varphi(z)\langle u,u\rangle=\overline{\varphi(u)}\langle u,z\rangle.
\]
A nonzero $\varphi$ cannot have $\varphi(u)=0$; thus all symmetric maps with that range have proportional functionals. For $n\ge2$ the coefficient ratios at $\sigma+it$ and $\sigma$ satisfy
\[
 \frac{v_n(\sigma+it)}{v_n(\sigma)}
 =\frac{\sigma+it}{\sigma}\,n^{-it}(1+O_{\sigma,t}(1/n)).
\]
Integers nearest $\exp(2\pi k/|t|)$ and $\exp((2\pi k+\pi)/|t|)$ give opposite nonzero subsequential limits. Hence no common positive metric symmetrizes this family for all critical heights. Restricting any proposed metric on the inserted family to $U\ell^2$ gives the same contradiction. An independently constructed positive comparison is still a possible different route.

\subsection{Precisely how classical and recent sources enter}
Over a finite field the Grothendieck trace-formula architecture writes a zeta function as alternating determinants on finite-dimensional cohomology; Deligne's purity supplies the weights of eigenvalues \cite{SGAFourHalf,DeligneWeilI,DeligneWeilII}. If a weight-$w$ Frobenius is $q^{w/2}$ times a positive-metric unitary, each eigenvalue has modulus $q^{w/2}$. The norm conclusion is immediate; constructing that form or proving purity is the deep step. Symplectic/alternating duality $F^tJF=q^wJ$ alone only pairs reciprocal eigenvalues and does not give their individual moduli. An order resolution or Schur-positive character expansion alone cannot replace this step in number-field cohomology.

The recent primary comparisons used here are specific:
\begin{itemize}
\item \cite[Theorem3.12, pp.15--17]{ConnesConsaniKnots} gives finite abelian covers with local orbit length $\log p$ and genuine arithmetic monodromy. It suggests what a descent-equipped order functor should match; it does not identify our $u$ or a face cycle with a prime without such a functor.
\item \cite[Propositions4.1--4.3, Theorem4.6]{Prolate} builds finite-prime local models with archimedean Fourier factors. Its Hilbertian isomorphisms are not an asserted uniformly controlled common isometry for our arithmetic approximants. \cite[Proposition3.1]{LocalMoments} gives determinacy of a particular local moment problem; positive finite Hankel matrices are not the missing global Weil inequality.
\item \cite[Theorem3.6]{ZetaTriples} gives discreteness/lower boundedness for the stated restricted Weil form. Theorem5.10, pp.23--24, additionally assumes a finite truncation, $\lambda>1$, a simple even bottom eigenfunction $\xi$ and $\delta_N\xi=1$. Its quotient metric is positive and its regularized determinant is $-i\lambda^{-iz}\widehat\xi(z)$. Passing to the infinite arithmetic comparison, including the ground-state and convergence conditions in its \S8, is not automatic.
\item \cite[Theorem6.1, p.13]{QuadraticForms} assumes a specified real distribution on an interval, the corresponding trigonometric-core form, essential self-adjointness, and an isolated simple lowest eigenvalue with reflection-even eigenfunction. Its Fourier transform then has only real zeros. Our degree operator's vacuum and length parity are not those data.
\item \cite[Theorem6.1]{DeningerPackets} gives prime-indexed periodic packets in its dynamical construction, not one proved orbit per prime with the full desired cohomological trace formula. \cite{DeningerReal} isolates an obstruction to a particular real-coefficient cohomology with full functorial Artin multiplicities; it is not a theorem excluding every Hilbert--P\'olya operator. The current survey \cite{ConnesSurvey} distinguishes remaining steps. The recent Picard-monoid construction \cite{ConnesConsaniJacobian} is not already a polarized abelian variety with a number-field purity theorem. \cite{PongeTian} uses RH as an input in the cited spectral identification, so cannot be used circularly as its proof.
\end{itemize}
These are targeted theorem/proof readings, not claims that every page of every cited work has been audited. The preserved review records specify actual page coverage and hashes, including inaccessible original PDFs where the modern primary proof was used instead.

\subsection{A concrete sequence of remaining constructions}
First choose a number-field model of a Belyi pair, its finite covers and conductor/ramification data. Construct actions or exact correspondences on the labelled order/module or analytic categories that recover the local prime Frobenius polynomials. The elliptic twists show why omitting this datum loses information. Test at more than one arithmetic model of the same complex dessin, not just its topological passport.

Second specify a finite projective or otherwise determinant-admissible complex and its \emph{independent} arithmetic endomorphism. Prove exact-sequence/determinant-line compatibility. Use semi-invariants to encode presentation degeneracies, and Schur/exterior characters to calculate local factors, retaining every central-action and nilpotence distinction above. General derived Schur replacements require their own construction; termwise Schur on the resolution fails the exhibited split sequence test.

Third prove the same local factors, gamma terms and conductor give a globally convergent or correctly regularized determinant. Trace-norm convergence must be established for the chosen analytic perturbations; strong convergence is insufficient in general. The inserted Abel family supplies this topology for an elementary scalar encoding, but has neither the independent arithmetic spectrum nor the positive metric needed for zero location.

Finally produce an independently proved positive form or self-adjoint determinant-preserving realization. One acceptance criterion is the Hurwitz target in the core; another is equality with the precise arithmetic quadratic form in the recent spectral theorems, with all their domain, ground-state, normalization and limit hypotheses verified. The local Cauchy kernel, $K$-group computation and graph spectral dimension alone meet none of these final comparison hypotheses. No remaining condition is silently asserted true here. Dirichlet, Hecke, Dedekind and Artin scopes must remain individually recorded.

\section{Graph pullbacks and completions with specified fibers}
\subsection{The complete Fock graph presentation}
For graph C* conventions the edge operator has initial projection at its range vertex and final projection under its source vertex. Put
\[
 S_h=L_{a_h}Q_{a_{\tau h}},\qquad T_h=L_{a_h}v_{[\tau h]}.
\]
The two right-hand projections are orthogonal subprojections of $L_{a_h}^*L_{a_h}=Q_{a_{\tau h}}+v_{[\tau h]}$. Thus $S_h^*S_h=Q_{a_{\tau h}}$, $T_h^*T_h=v_{[\tau h]}$, and $S_h^*T_h=0$. For distinct $h$ the prefix ranges are orthogonal, so every other cross-edge CK1 product is zero. Their final projections sum to $Q_{a_h}$. The vacua have no outgoing edges. The resulting graph has $|H|+|Q_0|$ vertices and $2|H|$ edges, with cycles only among its state vertices and exits to the paired sinks. Gauge covariance and nonzero vertices prove faithfulness as well as surjectivity. This explains precisely how a graph C* algebra is obtained without imposing the incompatible original-quiver CK1 relations.

The graph integer map has one column for each regular vertex $h$ and no sink columns:
\[
 d:\Z^H\longrightarrow\Z^H\oplus\Z^{Q_0},\qquad
 d(e_h)=e_h-e_{\tau h}-b_{[\tau h]}.
\]
The state part of $d(x)=0$ says $x_h=x_{\tau^{-1}h}$; the sink part then says the constant cycle coefficients are in $\ker B_\Gamma$. For the cokernel choose a base state on each cycle and successively use $q_h=q_{\tau h}+b_{[\tau h]}$ to eliminate all other states. Closing the cycle leaves $\sum_{h\in c}b_{[h]}=0$, while its base state remains free. Hence
\[
 \operatorname{coker}d\cong\Z^{\#\operatorname{cyc}\sigma}\oplus\operatorname{coker}B_\Gamma.
\]
This gives an independent integral check of the previous compact-extension computation. The full labelled graph, diagonal projections and edge families carry more information than those coarse groups.

\subsection{A genuine application of the supplied pullback theorem}
For the interval dessin let $E_2=G_\alpha$ have loop vertices $q_x,q_y$, common sink $b$, loops $s_x,s_y$ and exits $t_x,t_y$. Remove $b$ to obtain $F_2$, two disjoint loop graphs. Its quotient is exactly $C(\mathbb T)^2$. Let $E_1$ have the same three vertices and edges $e_{x,n}:q_x\to b$, $e_{y,n}:q_y\to b$, $n\ge0$, and no other edges. Remove $b$ to obtain the edgeless $F_1$, whose algebra is $\C^2$.

The inclusions $F_i\subset E_i$ are admissible in the precise sense of \cite[Definition3.1]{GraphPullbacks}: the removed sink set is hereditary/saturated; all edges ending at retained vertices are retained; no vertex emits infinitely many edges into the removed set and a positive finite number outside it. For $E_1$ the latter outside count is zero, which is allowed. Define a path functor identity on objects by
\[
 e_{x,n}\longmapsto s_x^nt_x,\qquad
 e_{y,n}\longmapsto s_y^nt_y.
\]
Images of distinct positive edges are prefix incomparable: at the first different length one has an exit and the other a loop. Source vertices and identities cause no exception. The image consists of all positive pointed paths in $E_2$, whose last edge is not a short loop; vertices are included separately. The first graph is acyclic and the removed sink has no short loop in the second. All hypotheses of Theorem3.4 are now verified, giving
\[
 C^*(E_1)\cong\mathcal T_A\mathop{\times}_{C(\mathbb T)^2}\C^2.
\]
The quotient of $C^*(E_1)$ by its sink ideal is $\C^2$, and its sink ideal is compact operators on the paths ending at $b$. The graph is acyclic, hence its algebra is AF. There are no regular vertices in $E_1$: the two sources are infinite emitters and $b$ is a sink. The graph $K$ matrix has no columns and three rows, yielding $K_0=\Z^3$ and $K_1=0$.

This is a positive instance of the new source inside our infinite dessin framework. It is a constants-symbol subalgebra of the Fock model, not the entire Fock algebra or the formal order. For a permitted cycle of length greater than one, its full closed path is itself pointed in this paper's convention, so an object-preserving acyclic path functor cannot cover all such paths. That simple obstruction prevents copying this example to every cycle by analogy. The sink-extension Theorem4.3 also requires a specified original pullback, injective gluing maps, annihilation of glued vertex projections and coverage of all paths ending at the gluing sinks. A two-to-one identification of sink occurrences within one branch does not satisfy these assumptions without another construction.

\subsection{Diagonal residue obstruction and the corrected pullback}
Suppose one tries to glue full matrix-circle blocks by their diagonal values. In $M_2$, $E_{12}E_{21}=E_{11}$, although both factors have zero diagonal. Thus the proposed boundary map is not a homomorphism. For $\ell>1$, more generally any *-homomorphism from the full matrix-circle algebra to a commutative C*-algebra has equal images of $E_{ii}=E_{ij}E_{ji}$ and $E_{jj}=E_{ji}E_{ij}$. These images are both orthogonal and equal, hence zero; summing the diagonal units kills the image of one and then the whole homomorphism. This works for nonunital *-maps as well: a nonzero image of one would give a nonzero unital corner map. The failure concerns full matrix-circle blocks, not the already-restricted fiber algebra $D_\ell$.

Evaluation $D_\ell\to\C^\ell$ is a genuine surjective *-homomorphism since at one both factors are diagonal. Its kernel is $C_0(\mathbb T\setminus\{1\},M_\ell)$. Pulling it back along the repeated-coordinate embedding $\Delta_\alpha:\C^{Q_0}\to\C^H$ gives exactly the fiber conditions of $\mathcal D_H$. The paired constant diagonal matrices give a unital *-section. Hence this is a split C*-extension, not a claim deduced from a merely commutative diagram of nonmultiplicative maps. Since $\mathbb T\setminus\{1\}\cong\R$, each kernel block is a matrix suspension of $\C$, with $K_0=0$, $K_1=\Z$. Split exactness now proves Theorem~\ref{thm:diagonalmodel}.

The pairing is retained in this specified presentation, not recovered from the unmarked isomorphism class. In particular arbitrary constant permutation matrices can relabel a branch's diagonal positions, so the spherical and elliptic theta examples have isomorphic unmarked $\mathcal D_H$ despite different ribbon orderings. Keep the labelled arrows, cycle cuts and distinguished fiber as additional data when using it as a dessin model. Also the central branch unitary $zI_\ell$ represents $\ell$ times the primitive corner generator in $K_1$, not a new independent generator per matrix entry.

\subsection{Generation, faithfulness and the exact failure of extension}
Let $\rho$ denote the polynomial representation in Theorem~\ref{thm:diagonalmodel}. Its vertex images are orthogonal projections summing to one. Every arrow image lies under the correct source and range vertex projections. A forbidden product vanishes because the intermediate state indices are $\tau h$ and $\alpha\tau h$, which are distinct. Thus $\rho$ is an actual algebra map, not a fitted characteristic polynomial.

For a full turn starting at $h$, $\rho(c_h)=w_cE_{hh}$ with $w_c(z)=g(z)^{|c|}z$. It is a normal element of $\mathcal D_H$. The function $w_c$ vanishes exactly at one; away from that point $w_c/|w_c|=z$. Therefore it is injective on all of $\mathbb T$, and a homeomorphism to its image. Nonunital continuous functional calculus at $\rho(c_h)$, using functions zero at zero, gives every $fE_{hh}$ with $f\in C(\mathbb T)$, $f(1)=0$. The zero condition makes the result vanish in all other matrix entries and branches.

To obtain an off-diagonal entry choose the unique shortest permitted path from state $h$ to $j$ within its cycle, of length $r>0$ and cut-crossing count $\kappa\in\{0,1\}$. Its image is $g^rz^\kappa E_{hj}$. Given $f(1)=0$, multiply it on the left by the already available diagonal element
\[
 \frac{f(z)z^{-\kappa}}{g(z)^r+\varepsilon}E_{hh},\qquad \varepsilon>0.
\]
The product is $f g^r/(g^r+\varepsilon)E_{hj}$ and converges uniformly to $fE_{hj}$. Indeed near one $|f|$ is small, and on the complement of that neighborhood $g^r$ has a positive minimum. These two bounds prove uniform convergence without assuming division by a function vanishing at the fiber. Thus the generated algebra contains the whole kernel ideal. Its quotient contains the entire paired diagonal algebra, so it is $\mathcal D_H$.

For injectivity let a finite linear combination of path images vanish. At one, all positive paths vanish, forcing the coefficients of the original vertex idempotents to vanish. Fix a state-matrix entry in a branch. Its remaining paths have lengths $r+m\ell$ and images
\[
 g^rz^\kappa\sum_m a_m(g^\ell z)^m E_{hj},
\]
with a possible initial zero coefficient if $r=0$. On $\mathbb T\setminus\{1\}$ the prefactor is nonzero, and $w_c$ has infinite range, so the polynomial has infinitely many distinct zeros. All coefficients are zero. Repeating by entries proves faithfulness.

No boundedness improvement repairs full formal extension. Pick any $z_0\ne1$, let $\lambda=w_c(z_0)^{-1}$, and consider $1-\lambda c_h$. Its inverse $\sum_{m\ge0}\lambda^m c_h^m$ exists in the arrow-adic order. The matrix image at $z_0$ has a zero eigenvalue on state $h$, so cannot be invertible. A unital algebra homomorphism would preserve the inverse identity, an immediate contradiction. This proves nonextension without any continuity assumption. It does not rely on the false claim that a norm greater than one alone forces a geometric-series denominator to be singular. Scaling $g$ to have norm at most one still does not allow arbitrary formal series.

The two valid operator completions are structurally different. In the bipartite case $K_0(\mathcal T_A)=\Z^{E+1}$ differs from $K_0(\mathcal D_H)=\Z^E$; in the nonbipartite case the former has the extra $\Z/2$ torsion. Neither is automatically the norm completion of the arrow-adic topology. No invariant constructed solely from these complex objects distinguishes the rational twists already exhibited; arithmetic descent remains a separate input.

\subsection{Invariant basis number and the Leavitt comparison}
If a unital $k$-algebra $R_1$ has a unital representation on a nonzero finite-dimensional vector space $V$, an isomorphism $R_1^m\cong R_1^n$ induces an isomorphism $V^m\cong V^n$ by tensoring. Comparing dimensions gives $m=n$. This proves invariant basis number without requiring the representation to be faithful. Every original dessin algebra and completed order has a vertex-simple representation, and each permitted-state Laurent algebra has evaluation onto a nonzero matrix block. The Leavitt algebra of $G_\alpha$ has the same finite-dimensional representation after killing the vacuum sinks and evaluating one boundary cycle. Thus all of these examples have invariant basis number. Leavitt's free-rank phenomena are important elsewhere but cannot be assumed in this class simply from the name LPA.

For the specified finite graph $G_\alpha$, the usual complex Leavitt path algebra embeds in $C^*(G_\alpha)$ with norm-dense image, using its canonical generators and involution \cite[Theorem7.3]{TomfordeEmbedding}. Those generators are the split $S_h,T_h$, not the unsplit original gentle arrows equipped with ordinary CK1. This is the exact algebraic/analytic twin relationship applicable here.

\section{String finiteness and the smooth quotient surface}
For a finite quiver, finitely controlled implies finitely generated: choose one finitely generated submodule containing each vertex space and sum the finitely many choices. Conversely the module itself works. Stabilization at each of finitely many vertices likewise gives ordinary artinianity. Crawley--Boevey's polynomial theorem therefore applies directly to our finitely generated and artinian polynomial modules, with its primitive injective bands in the latter case. It does not assert a classification of every possible algebraic module.

The formal branch $M=k[[x]]$, $y=0$, gives a precise counterexample to that stronger assertion. An $A$-endomorphism $F$ preserves each $x^nM$. For the polynomial truncation $f_n$ of $f$, one has $F(f)-f_nF(1)\in x^nM$. Separatedness and all $n$ imply $F(f)=fF(1)$. Hence $\operatorname{End}_A(M)=k[[x]]$ is local, so $M$ is indecomposable. If an ordinary word module has $y=0$, its word cannot contain $y$ or $y^{-1}$, since such a letter would create a nonzero basis map. A reduced word using only $x^{\pm1}$ cannot change direction. Its nonperiodic choices are finite strings, the free ray $k[x]$, and the opposite locally nilpotent ray. The first and last fail injectivity of $x$, and the free ray fails invertibility of $1-x$. A two-sided pure-$x$ word is periodic, and in every associated band $x$ acts invertibly. On $M$, $x$ is injective but not onto. None of these possibilities equals $M$. This includes arbitrary nonzero indecomposable Laurent coefficient modules, not just finite bands.

For a general nonzero completed lattice, finite generation over $A$ would imply finite generation over $k[u]$. Every nonzero polynomial in $u$ is a power of $u$ times a unit of $k[[u]]$, so it acts injectively. The lattice would be a nonzero finite free $k[u]$-module. Multiplication by $1-u$ on such a free module has a nonzero cokernel, but the same map is invertible over $k[[u]]$. This proves the contradiction in the core. The strictly descending $u$-adic chain separately excludes artinianity. These are actual failed hypotheses, not merely unverified applicability.

Lekili--Polishchuk's precise convention takes $D(A)$ to mean perfect dg modules. Their locally gentle definition allows infinite dimension, and their smoothness criterion says no forbidden cycle; properness says no permitted cycle. Both fail for every nonempty closed dessin here. Their Theorem3.2.2 therefore does not apply to the full polynomial dessin algebra, even though it applies to suitable infinite-dimensional quotients. The OPS finite-dimensional geometric-intersection theorem has its own grading, local-system, boundary and exceptional coincident-curve conditions. An unrestricted statement that every gentle/string derived object is an ungraded geodesic crossing is not justified by either source.

For the interval quotient $B=A/(y)=k[x]$, global dimension one follows from its PID structure: every submodule of a finite free module is free, giving finite projective resolutions of all finitely generated modules. Thus $\operatorname{Perf}(B)=D^b(\operatorname{mod}_{\rm fg}B)$ in grading zero. Its permitted loop gives one unmarked boundary component and its length-one forbidden thread $x$ gives one stop on the other boundary, the annulus of the smooth partially wrapped model. This construction uses the infinite polynomial quotient, not the completed $k[[x]]$ category.

Finally the quotient does not give a derived fully faithful embedding by ordinary inflation. Resolve $B$ over $A$ by the alternating maps $y,x,y,x,\ldots$. Hom into $B$ gives $0,x,0,x,\ldots$, so $\operatorname{Ext}_A^{2r}(B,B)=B/xB=k$ for every $r\ge1$, whereas all positive self-Ext over $B$ vanish. Consequently its inflation is not derived fully faithful, and the perfect $B$-module $B$ becomes an infinite-projective-dimension $A$-module. A surface functor transferring these categories must account for that difference explicitly.


\section{Literal formal and computational coverage}
\texttt{DessinNodeOrder.lean} defines the fiber-product subring of two genuine power-series rings. Its twenty printed declarations establish diagonal injectivity, nonzero branches, branch zero products, factorizations and annihilator identities, regular $u$, augmentation surjectivity and exactness, bijective power-series coordinates, square-zero alternating maps, three kernel-equals-range statements and projectivity of the actual regular/product terms. The coordinate statement is a bijection; the file does not yet package it as an $R$-linear equivalence or declare the general order freeness theorem.

\texttt{DessinInvariantComponents.lean} has ten printed declarations: three CK1-collapse components, determinant change of bases with an actual inverse, symmetric-square zero, identity, composition and relation preservation, the degree-two Schur character identity, and the nodal central-parameter obstruction. Its symmetric-square matrix is
\[
 \operatorname{Sym}^2\!\begin{pmatrix}a&b\\c&d\end{pmatrix}
 =\begin{pmatrix}a^2&ab&b^2\\2ac&ad+bc&2bd\\c^2&cd&d^2\end{pmatrix}.
\]
This is the ordinary monomial basis convention; composition is proved universally over any commutative ring. The full project suite prints types and foundational dependencies; no theorem about general LR decompositions, abstract derived equivalences, C* $K$-theory or GRH is inferred from these component proofs.

The exact Python checks use fixed rational matrices and exhaustive finite sets, with no random seeds. They construct the displayed permutations and paths, evaluate rational Cartan identities, check formal Euler products through a stated degree, verify unrestricted-output Fock basis identities on bounded input sets, compute incidence Smith data, construct all nine string representations and signed maps, and test Young/LR/Cauchy examples and the seven-element-field counts. Those diagnostics are reproducible finite corroboration. The universal infinite claims rest on the proofs or explicitly cited theorems, and the arithmetic comparison remains open.

The ten additional \texttt{DessinPullbackComponents.lean} statements use literal two-state matrix units and a subring of a product of arbitrary rings, with equal images under two ring homomorphisms. The section hypothesis of its projection-surjectivity theorem is explicit. It does not construct a topological C*-pullback or assert that a commutative diagram is automatically a pullback. The no-residue theorem is universal algebra, with a nontrivial commutative target; it is not an assumed arithmetic positivity premise.
